\ssr{ЗАКЛЮЧЕНИЕ}  

В ходе выполнения курсовой работы разработано программное решение для генерации и трёхмерной визуализации городской среды, полностью удовлетворяющее поставленному заданию. Все задачи, сформулированные во введении, успешно реализованы:  
\begin{itemize}  
    \item создан алгоритм генерации транспортной сети с тремя и более перекрёстками, гарантирующий отсутствие изолированных участков;
    \item созданы объемные модели зданий пяти архитектурных типов с минимальной детализацией 20 полигонов на объект и алгоритм их размещения вдоль дорог без взаимных пересечений;
    \item обеспечена воспроизводимость генерации за счёт фиксации начального значения генератора псевдослучайных чисел;
    \item реализовано интерактивное управление камерой, а именно, перемещение, вращение, масштабирование;
    \item интегрирован алгоритм карты теней для расчёта динамического освещения с возможностью настройки источников света.
\end{itemize}  

Анализ результатов, полученных в ходе исследования, позволяет сделать следующие выводы:  
\begin{itemize}  
    \item в аналитической части разобраны математические модели аффинных преобразований, определены методы, использующиеся в программе, а именно Z-буфер для определения перекрытия полигонов, диффузной закраски и карты теней для решения задачи теней и освещения;
    \item в конструкторской части, разработан алгоритм генерации городской среды и отрисовки кадра, а также предложен подход к распараллеливанию растеризации, что повысило производительность обработки сложных сцен;
    \item реализация на базе C++20 и Qt 6.5 позволила создать модульную архитектуру приложения с чётким разделением логики генерации, рендеринга и управления пользовательским интерфейсом;
    \item исследования показали квадратичную зависимость времени отрисовки от разрешения кадра и выявили константную задержку при обработке большого числа полигонов, что обусловлено необходимостью перебора всех объектов сцены.
\end{itemize}

Таким образом, цель работы --- создание программного комплекса для автоматической генерации и интерактивной визуализации городской среды с использованием современных методов рендеринга --- достигнута. Результаты исследования демонстрируют возможность применения процедурной генерации в связке с алгоритмами динамического освещения для создания масштабируемых 3D-сцен, а полученные данные о производительности предоставляют основу для дальнейшей оптимизации решения.  

\clearpage